% CVPR 2026 Paper Template; see https://github.com/cvpr-org/author-kit

\documentclass[10pt,twocolumn,letterpaper]{article}

%%%%%%%%% PAPER TYPE  - PLEASE UPDATE FOR FINAL VERSION
% \usepackage{cvpr}              % To produce the CAMERA-READY version
\usepackage[review]{cvpr}      % To produce the REVIEW version
% \usepackage[pagenumbers]{cvpr} % To force page numbers, e.g. for an arXiv version

% Import additional packages in the preamble file, before hyperref
% Additional packages for the paper

% Math
\usepackage{amsmath}
\usepackage{amssymb}
\usepackage{bm}

% Tables
\usepackage{booktabs}
\usepackage{multirow}
\usepackage{makecell}
\usepackage{colortbl}
\usepackage{array}

% Figures
\usepackage{graphicx}
\usepackage{subcaption}
\usepackage{float}
\usepackage{placeins}

% Algorithms
\usepackage{algorithm}
\usepackage{algpseudocode}

% Misc
\usepackage{xspace}
\usepackage{enumitem}
\usepackage{xcolor}
\usepackage{microtype}
\usepackage{soul}
\usepackage{pifont}

% Checkmark and crossmark
\newcommand{\cmark}{\ding{51}}
\newcommand{\xmark}{\ding{55}}

% Custom commands (etal, ie, eg, cf, wrt already defined by cvpr.sty)
\providecommand{\etal}{\textit{et al.}\xspace}
\providecommand{\ie}{\textit{i.e.}\xspace}
\providecommand{\eg}{\textit{e.g.}\xspace}
\providecommand{\cf}{\textit{cf.}\xspace}
\providecommand{\wrt}{w.r.t.\xspace}
\newcommand{\dataset}{\textsc{RealDeg-Bench}\xspace}
\newcommand{\pipeline}{\textsc{PhysDeg}\xspace}

\DeclareMathOperator*{\argmin}{arg\,min}
\DeclareMathOperator*{\argmax}{arg\,max}

% Table highlight colors
\definecolor{bestcolor}{HTML}{E8F5E9}
\definecolor{secondcolor}{HTML}{FFF3E0}

% Annotations (disable for submission)
\newcommand{\red}[1]{{\color{red}#1}}
\newcommand{\todo}[1]{{\color{red}#1}}
\newcommand{\TODO}[1]{\textbf{\color{red}[TODO: #1]}}


% It is strongly recommended to use hyperref, especially for the review version.
% hyperref with option pagebackref eases the reviewers' job.
% Please disable hyperref *only* if you encounter grave issues, 
% e.g. with the file validation for the camera-ready version.
%
% If you comment hyperref and then uncomment it, you should delete *.aux before re-running LaTeX.
% (Or just hit 'q' on the first LaTeX run, let it finish, and you should be clear).
\definecolor{cvprblue}{rgb}{0.21,0.49,0.74}
\usepackage[pagebackref,breaklinks,colorlinks,allcolors=cvprblue]{hyperref}
\usepackage[capitalize,noabbrev]{cleveref}

%%%%%%%%% PAPER ID  - PLEASE UPDATE
\def\paperID{*****} % *** Enter the Paper ID here
\def\confName{CVPR}
\def\confYear{2026}

%%%%%%%%% TITLE - PLEASE UPDATE
\title{Beyond Bicubic: Benchmarking Single Image Super-Resolution Under Physically-Grounded Complex Degradations}

%%%%%%%%% AUTHORS - PLEASE UPDATE
\author{Anonymous CVPR submission\\
Paper ID \paperID}

\begin{document}
\maketitle
% =================================================================================================
% ABSTRACT
% =================================================================================================
\begin{abstract}
\noindent
Single Image Super-Resolution (SISR) has witnessed remarkable progress, yet the dominant evaluation paradigm relies on benchmarks with simple bicubic downsampling, failing to capture the complex degradations in real-world imaging.
This mismatch leads to inflated metrics that poorly predict in-the-wild generalization.

We present \dataset, a large-scale benchmark accompanied by \pipeline, a physically-grounded, multi-stage degradation pipeline modeling the image formation process from scene radiance to stored pixel values across six stages: optical aberrations, motion blur, scene illumination, sensor noise, ISP artifacts, and lossy compression.
We construct \dataset by applying \pipeline to 2{,}000 high-quality images spanning four semantic categories, yielding paired LR--HR data with comprehensive per-image degradation metadata.

Evaluation of four state-of-the-art methods reveals significant performance degradation: models achieving 30+ dB PSNR on Set5 attain only $\sim$23~dB on \dataset, exposing a 6--8~dB generalization gap.
Furthermore, models designed for blind SR (\eg, Real-ESRGAN) exhibit superior perceptual quality despite lower distortion metrics, underscoring the need for multi-faceted evaluation.
Our benchmark and code will be publicly released.
\end{abstract}

% =================================================================================================
% 1. INTRODUCTION
% =================================================================================================
\section{Introduction}
\label{sec:intro}

Single Image Super-Resolution (SISR) aims to reconstruct a high-resolution (HR) image $I_{HR} \in \mathbb{R}^{sH \times sW \times 3}$ from its low-resolution (LR) counterpart $I_{LR} \in \mathbb{R}^{H \times W \times 3}$, where $s$ denotes the upsampling scale factor.
This ill-posed inverse problem has broad applicability in medical imaging~\cite{shi2013cardiac}, remote sensing~\cite{shermeyer2019effects}, surveillance~\cite{zhang2010super}, and consumer photography.
The field has progressed from dictionary-based methods~\cite{yang2010image} through CNNs~\cite{dong2015image,lim2017enhanced} to Vision Transformers~\cite{liang2021swinir,chen2023hat} and GANs~\cite{wang2021real,wang2018esrgan}, with top models exceeding 32~dB PSNR on Set5~\cite{bevilacqua2012low} at $\times$4.

However, the dominant benchmarking paradigm generates LR images through a single bicubic downsampling: $I_{LR} = I_{HR} \downarrow_{s}^{\text{bicubic}}$.
This assumes an idealized pipeline where the sole information loss is spatial resolution reduction via a known kernel.
In practice, real cameras introduce optical aberrations, sensor noise, motion blur, ISP artifacts, and lossy compression~\cite{liu2022blind,zhang2021designing}.
Models trained under the bicubic assumption often fail catastrophically on these real-world degradations~\cite{wang2021real}, a phenomenon we term the \textit{degradation gap}.

\begin{figure*}[t]
    \centering
    \includegraphics[width=\textwidth]{figures/fig-1.png}
    \caption{\textbf{Overview of the \pipeline degradation pipeline.}
    Six sequential stages model the image formation process from scene to stored file:
    (1)~optical aberrations,
    (2)~motion blur (5 kernel types),
    (3)~scene illumination effects,
    (4)~sensor noise,
    (5)~ISP processing, and
    (6)~compression and downsampling.
    Each stage is parameterized by physically meaningful quantities and applied stochastically, yielding diverse LR--HR pairs with complete degradation metadata.}
    \label{fig:pipeline_overview}
\end{figure*}

In this paper, we address this gap through two contributions.
First, we introduce \pipeline, a \textit{physically-grounded, multi-stage degradation pipeline} organizing ten distinct degradation operations into six ordered stages---optical, motion, scene, sensor, ISP, and compression---each parameterized by quantities with clear physical interpretations.
Second, we present \dataset, a large-scale benchmark of 2{,}000 diverse image pairs across four semantic categories (\textit{urban architecture}, \textit{natural landscapes}, \textit{portraits}, \textit{macro objects}) with comprehensive per-image degradation metadata.

Our evaluation of four SOTA methods reveals:
\begin{enumerate}[leftmargin=*,itemsep=1pt,topsep=2pt]
    \item \textbf{Substantial generalization gap:} models achieving ${>}$30~dB on Set5 attain only ${\sim}$23--24~dB on \dataset (6--8~dB drop).
    \item \textbf{Perception--distortion divergence:} Real-ESRGAN obtains the best LPIPS (0.444 vs.\ ${>}$0.55 for non-blind methods) despite the lowest PSNR, showing that distortion metrics alone are insufficient under complex degradations.
    \item \textbf{Content-dependent vulnerability:} urban architecture is most challenging due to geometric-degradation interactions; macro objects are comparatively resilient.
\end{enumerate}

\noindent Our contributions can be summarized as:
\begin{itemize}[leftmargin=*,itemsep=1pt,topsep=2pt]
    \item We propose \pipeline, a physically-grounded degradation pipeline comprising six ordered stages with ten degradation types, each governed by physically meaningful parameters.
    \item We construct \dataset, a benchmark of 2{,}000 image pairs spanning four semantic categories with complete per-image degradation metadata.
    \item We quantify the degradation gap: SOTA methods suffer a 6--8~dB PSNR drop from bicubic to realistic degradations, motivating future research in degradation-robust SR.
\end{itemize}

% =================================================================================================
% 2. RELATED WORK
% =================================================================================================
\section{Related Work}
\label{sec:related}

\paragraph{Single image super-resolution.}
Deep SISR has evolved rapidly from SRCNN~\cite{dong2015image} through residual architectures (EDSR~\cite{lim2017enhanced}, RCAN~\cite{zhang2018image}) to Transformer-based methods.
SwinIR~\cite{liang2021swinir} adapted shifted-window attention for efficient image restoration, and subsequent works---HAT~\cite{chen2023hat}, DSCF-SR~\cite{xiao2024dscf}, SeemoRe~\cite{zamfir2024seemore}---have further pushed quantitative benchmarks.
In parallel, GAN-based approaches (SRGAN~\cite{ledig2017photo}, ESRGAN~\cite{wang2018esrgan}) demonstrated that perceptual and adversarial losses produce photo-realistic textures, establishing a fundamental tension between distortion and perceptual quality~\cite{blau2018perception}.
Recent perceptual objectives include self-supervised DINO features~\cite{caron2021emerging} and frequency-domain losses~\cite{fuoli2021fourier}.

\paragraph{Degradation modeling.}
The classical degradation model $I_{LR} = (I_{HR} \otimes k) \downarrow_s + n$~\cite{dong2015image} is typically simplified to bicubic downsampling with $n{=}0$.
Blind SR methods address more complex scenarios: IKC~\cite{gu2019blind} iteratively corrects kernels, BSRGAN~\cite{zhang2021designing} randomly shuffles blur-downsample-noise operations, and Real-ESRGAN~\cite{wang2021real} applies a second-order degradation pipeline with sinc filters.
Despite advances in degradation-aware \textit{training}, the \textit{evaluation} paradigm has not kept pace---models are still predominantly assessed on bicubically degraded benchmarks.
Our work draws on camera ISP modeling~\cite{zamir2020cycleisp} and sensor noise characterization~\cite{healey1994radiometric,foi2008practical} to construct a physically principled pipeline.

\paragraph{SISR benchmark datasets.}
The widely used benchmarks---Set5~\cite{bevilacqua2012low}, Set14~\cite{zeyde2010single}, BSD100~\cite{martin2001database}, Urban100~\cite{huang2015single}---all employ bicubic downsampling.
Real-world alternatives like RealSR~\cite{cai2019toward} (595 pairs, single camera) and DRealSR~\cite{wei2020component} (840 pairs, five cameras) provide genuine degradations but are limited in scale and diversity, with potential alignment artifacts.
Our \dataset occupies a distinct niche: unlike bicubic benchmarks, we model diverse, physically-grounded degradations; unlike real-captured datasets, our synthetic approach provides exact ground truth without alignment errors.
Uniquely, \dataset includes per-image degradation metadata enabling fine-grained robustness analysis (\Cref{tab:dataset_comparison}).

\begin{table}[t]
\centering
\caption{\textbf{Comparison of SISR benchmark datasets.} \dataset uniquely combines large scale, diverse categories, complex degradations, and per-image metadata.}
\label{tab:dataset_comparison}
\resizebox{\columnwidth}{!}{%
\begin{tabular}{@{}lcccccc@{}}
\toprule
\textbf{Dataset} & \textbf{Size} & \textbf{Degrad.} & \textbf{Categories} & \textbf{Metadata} & \textbf{GT} \\
\midrule
Set5~\cite{bevilacqua2012low} & 5 & Bicubic & \xmark & \xmark & \cmark \\
Set14~\cite{zeyde2010single} & 14 & Bicubic & \xmark & \xmark & \cmark \\
BSD100~\cite{martin2001database} & 100 & Bicubic & \xmark & \xmark & \cmark \\
Urban100~\cite{huang2015single} & 100 & Bicubic & 1 & \xmark & \cmark \\
DIV2K~\cite{agustsson2017ntire} & 1{,}000 & Bicubic & \xmark & \xmark & \cmark \\
RealSR~\cite{cai2019toward} & 595 & Real (1 cam.) & \xmark & \xmark & $\sim$ \\
DRealSR~\cite{wei2020component} & 840 & Real (5 cam.) & \xmark & \xmark & $\sim$ \\
\midrule
\textbf{\dataset} & \textbf{2{,}000} & \textbf{Complex} & \textbf{4} & \textbf{\cmark} & \textbf{\cmark} \\
\bottomrule
\end{tabular}%
}
\end{table}

% =================================================================================================
% 3. PROPOSED DEGRADATION PIPELINE
% =================================================================================================
\section{Proposed Degradation Pipeline}
\label{sec:pipeline}

We introduce \pipeline, a physically-grounded degradation pipeline that models the sequential transformation of scene radiance into a stored low-resolution image.
Unlike existing approaches that apply degradations in a flat or shuffled manner~\cite{zhang2021designing,wang2021real}, our pipeline organizes degradation operations into six physically-ordered stages that mirror the actual image formation process:
\begin{equation}
    I_{LR} = \left( \mathcal{D}_{\text{comp}} \circ \mathcal{D}_{\text{isp}} \circ \mathcal{D}_{\text{sensor}} \circ \mathcal{D}_{\text{scene}} \circ \mathcal{D}_{\text{motion}} \circ \mathcal{D}_{\text{optical}} \right)(I_{HR}) \downarrow_{s},
    \label{eq:full_pipeline}
\end{equation}
where each operator $\mathcal{D}_{(\cdot)}$ is parameterized by physically meaningful quantities and applied stochastically according to a difficulty-modulated probability.

\subsection{Stage 1: Optical Degradations}
\label{sec:pipeline_optical}

Optical degradations originate from imperfections in the lens system.
We model three effects: \textit{chromatic aberration} via wavelength-dependent channel shifts $\delta_c \sim \mathcal{U}(-2.5, 2.5)$ pixels~\cite{kang2007automatic};
\textit{radial lens distortion} using the Brown--Conrady model~\cite{brown1971close}:
\begin{equation}
    r' = r \left(1 + k_1 r^2 \right), \quad k_1 \sim \mathcal{U}(-0.15, 0.15);
    \label{eq:radial_distortion}
\end{equation}
and \textit{vignetting} via the $\cos^4$ fall-off approximation $V(x,y) = 1 - s \cdot (d(x,y)/d_{\max})^2$ with $s \sim \mathcal{U}(0.1, 0.5)$.

\subsection{Stage 2: Motion Blur}
\label{sec:pipeline_motion}

Motion blur arises during the exposure interval and is modeled as $I_{\text{blur}} = I \ast k_\theta$.
We implement five physically-motivated kernel types to capture the diversity of real-world motion:
(1)~\textit{linear} uniform kernels along direction $\phi$ with length $l \sim \mathcal{U}(5, 20)$~px;
(2)~\textit{defocus} disk kernels of radius $r \sim \mathcal{U}(2, 10)$~px, approximating circular-aperture bokeh~\cite{potmesil1981lens};
(3)~\textit{anisotropic Gaussian} kernels with independent scales $\sigma_x, \sigma_y \sim \mathcal{U}(0.5, 4.0)$ and rotation $\phi \sim \mathcal{U}(0, \pi)$;
(4)~\textit{rotational} blur simulating in-plane camera rotation; and
(5)~\textit{random walk} trajectories $\mathbf{p}_t = \sum_{\tau=1}^{t} \boldsymbol{\epsilon}_\tau$ with $\boldsymbol{\epsilon}_\tau \sim \mathcal{N}(\mathbf{0}, \sigma_{\text{step}}^2 \mathbf{I}_2)$ simulating camera shake.

\subsection{Stage 3: Scene Illumination Effects}
\label{sec:pipeline_scene}

We model two scene-level effects.
\textit{Low-light degradation} applies gamma correction with noise amplification: $I_{\text{low}} = I^\gamma + \mathcal{N}(0, \sigma_0^2 b^2)$ where $\gamma \sim \mathcal{U}(1.5, 3.0)$ and $b \sim \mathcal{U}(1.2, 2.0)$, modeling the increased noise visibility in underexposed imagery~\cite{chen2018learning}.
\textit{Atmospheric scattering} follows the Koschmieder model~\cite{mccartney1976optics}: $I_{\text{haze}} = t \cdot J + A(1 - t)$, where $A \sim \mathcal{U}(0.8, 0.95)$ is the atmospheric light and $t = \exp(-\beta \cdot d)$ with $\beta \sim \mathcal{U}(0.1, 0.6)$.

\subsection{Stage 4: Sensor Noise}
\label{sec:pipeline_sensor}

Camera sensor noise is a composite of multiple physical sources~\cite{healey1994radiometric,foi2008practical}.
We model three components: (1)~signal-dependent \textit{shot noise} (Poisson):
\begin{equation}
    I_{\text{shot}}(x, c) = \text{Poisson}\!\left(\frac{I(x, c)}{\alpha_c}\right) \cdot \alpha_c,
    \label{eq:shot_noise}
\end{equation}
where $\alpha_c = \alpha \cdot \lambda_c$ with $\alpha \sim \mathcal{U}(0.01, 0.08)$ and channel factors $[\lambda_R, \lambda_G, \lambda_B] = [1.1, 1.0, 1.05]$;
(2)~spatially-correlated \textit{read noise} (Gaussian) with $\sigma_r \sim \mathcal{U}(0.005, 0.025)$; and
(3)~\textit{dark current} offset $d \sim \mathcal{U}(0.001, 0.012)$.
The composite model is $I_{\text{noisy}} = \text{clip}[\text{Shot}(I) + \text{Read}(I) + d,\, 0,\, 1]$.

\subsection{Stage 5: Image Signal Processing (ISP)}
\label{sec:pipeline_isp}

Consumer cameras apply an ISP pipeline that introduces its own artifacts.
We model five sub-stages in their canonical processing order: $\mathcal{D}_{\text{isp}} = \mathcal{S} \circ \mathcal{T} \circ \mathcal{C} \circ \mathcal{F} \circ \mathcal{M}$, where $\mathcal{M}$ denotes demosaicing (cross-channel leakage simulation), $\mathcal{F}$ bilateral denoising~\cite{tomasi1998bilateral}, $\mathcal{C}$ stochastic color correction via a perturbed $3{\times}3$ matrix with saturation adjustment, $\mathcal{T}$ gamma-based tone mapping with highlight compression, and $\mathcal{S}$ unsharp masking.

\subsection{Stage 6: Compression and Downsampling}
\label{sec:pipeline_compression}

The final stage applies JPEG compression ($q \sim \mathcal{U}(25, 85)$), an optional random intermediate resize ($f \sim \mathcal{U}(0.7, 1.3)$) simulating social media processing, and terminal downsampling by factor $s$ using a randomly selected kernel from $\{\text{bicubic, bilinear, Lanczos}\}$.

\subsection{Difficulty Modulation}
\label{sec:pipeline_difficulty}

Each degradation is applied independently with a difficulty-modulated probability.
We define four levels---\textit{easy}, \textit{medium}, \textit{hard}, \textit{extreme}---controlling the base application probability $p$ and maximum concurrent degradations $N_{\max}$:
$(p, N_{\max}) \in \{(0.3, 2), (0.5, 3), (0.7, 5), (0.9, 7)\}$.
Stage-specific modifiers further refine rates: optical effects are scaled by $0.7{\times}$ (less common), while ISP and sensor noise are scaled by $1.2{\times}$ and $1.1{\times}$ (highly prevalent).
\Cref{tab:degradation_params} summarizes all degradation types and their parameters.

\begin{table}[t]
\centering
\caption{\textbf{Summary of degradation types in \pipeline.} Application probabilities shown for the default random difficulty mode.}
\label{tab:degradation_params}
\resizebox{\columnwidth}{!}{%
\begin{tabular}{@{}llcc@{}}
\toprule
\textbf{Stage} & \textbf{Degradation} & \textbf{Key Parameters} & \textbf{Prob.} \\
\midrule
\multirow{3}{*}{\rotatebox{90}{\small Optical}}
& Chromatic aberr. & $\delta \in [-2.5, 2.5]$ px & $\sim$0.4 \\
& Lens distortion & $k_1 \in [-0.15, 0.15]$ & $\sim$0.4 \\
& Vignetting & $s \in [0.1, 0.5]$ & $\sim$0.3 \\
\midrule
\rotatebox{90}{\small Motion}
& Motion blur & 5 kernel types & 0.50 \\
\midrule
\multirow{2}{*}{\rotatebox{90}{\small Scene}}
& Low-light & $\gamma \in [1.5, 3.0]$ & $\sim$0.4 \\
& Atmospheric & $\beta \in [0.1, 0.6]$ & $\sim$0.3 \\
\midrule
\rotatebox{90}{\small Sensor}
& Composite noise & $\alpha, \sigma_r, d$ & $\sim$0.9 \\
\midrule
\rotatebox{90}{\small ISP}
& 5-stage ISP & See \cref{sec:pipeline_isp} & $\sim$0.8 \\
\midrule
\multirow{2}{*}{\rotatebox{90}{\small Comp.}}
& JPEG compress. & $q \in [25, 85]$ & 0.70 \\
& Random resize & $f \in [0.7, 1.3]$ & 0.60 \\
\bottomrule
\end{tabular}%
}
\end{table}

% =================================================================================================
% 4. DATASET CONSTRUCTION
% =================================================================================================
\section{\dataset: Dataset Construction}
\label{sec:dataset}

\subsection{Source Image Collection}
\label{sec:dataset_collection}

We curate 2{,}000 high-quality HR images from Pexels\footnote{\url{https://www.pexels.com}} (CC0 license), organized into four semantically distinct categories:
\textbf{Urban Architecture} (${\sim}$500): buildings and cityscapes with geometric patterns and fine structural details;
\textbf{Natural Landscapes} (${\sim}$500): outdoor scenes with organic textures and complex foliage;
\textbf{Portraits} (${\sim}$500): human subjects with smooth skin gradients and fine hair details;
\textbf{Macro Objects} (${\sim}$500): close-up photography with fine-grained textures.
All sources satisfy a minimum resolution of $1024 \times 1024$ pixels.

\subsection{Generation Process}
\label{sec:dataset_generation}

\Cref{alg:dataset_generation} outlines the procedure.
For each HR image, we apply the \pipeline pipeline (\cref{sec:pipeline}) in random difficulty mode ($p \sim \mathcal{U}(0.4, 0.8)$, $N_{\max} \sim \mathcal{U}(2, 6)$), then downsample by $s{=}4$.
Complete degradation metadata $\mathbf{m}_i$ is recorded for every pair---the exact sequence of applied degradations, their parameters, and difficulty configuration---enabling stratified evaluation, degradation-conditioned model training, and detailed failure analysis.

\begin{algorithm}[t]
\caption{Dataset generation procedure}
\label{alg:dataset_generation}
\begin{algorithmic}[1]
\Require Source images $\{I_i^{HR}\}_{i=1}^{N}$, scale factor $s$, pipeline $\mathcal{P}$
\Ensure Paired dataset $\{(I_i^{LR}, I_i^{HR}, \mathbf{m}_i)\}_{i=1}^{N}$
\For{$i = 1$ \textbf{to} $N$}
    \State Sample difficulty: $p \sim \mathcal{U}(0.4, 0.8)$, $N_{\max} \sim \mathcal{U}(2, 6)$
    \State Initialize metadata: $\mathbf{m}_i \leftarrow \emptyset$
    \For{each stage $\mathcal{D}_k$ in $\mathcal{P}$}
        \State Sample application: $u \sim \mathcal{U}(0, 1)$
        \If{$u < p \cdot w_k$ \textbf{and} $|\mathbf{m}_i| < N_{\max}$}
            \State $I_i^{\text{deg}}, \theta_k \leftarrow \mathcal{D}_k(I_i^{\text{deg}}; \text{random})$
            \State $\mathbf{m}_i \leftarrow \mathbf{m}_i \cup \{(k, \theta_k)\}$
        \EndIf
    \EndFor
    \State $I_i^{LR} \leftarrow \text{JPEG}(I_i^{\text{deg}};\, q) \downarrow_{\phi, s}$
\EndFor
\end{algorithmic}
\end{algorithm}

\subsection{Dataset Statistics}
\label{sec:dataset_stats}

\begin{table}[t]
\centering
\caption{\textbf{Dataset statistics of \dataset.}}
\label{tab:dataset_stats}
\resizebox{\columnwidth}{!}{%
\begin{tabular}{@{}lcccc@{}}
\toprule
\textbf{Statistic} & \textbf{Urban} & \textbf{Landscape} & \textbf{Portrait} & \textbf{Macro} \\
\midrule
Number of pairs & $\sim$500 & $\sim$500 & $\sim$500 & $\sim$500 \\
Min HR resolution & \multicolumn{4}{c}{$1024 \times 1024$ pixels} \\
Scale factor & \multicolumn{4}{c}{$\times 4$} \\
\midrule
\multicolumn{5}{l}{\textit{Degradation coverage (mean applied per image):}} \\
Optical degradations & \multicolumn{4}{c}{0.8--1.2 per image} \\
Motion blur & \multicolumn{4}{c}{0.3--0.5 per image} \\
Sensor noise & \multicolumn{4}{c}{0.7--0.9 per image} \\
ISP artifacts & \multicolumn{4}{c}{2.5--4.0 stages per image} \\
JPEG compression & \multicolumn{4}{c}{0.65--0.75 per image} \\
\midrule
Avg. degradations/image & \multicolumn{4}{c}{3.2 $\pm$ 1.4} \\
Difficulty range & \multicolumn{4}{c}{Easy -- Extreme (continuous)} \\
Metadata fields/image & \multicolumn{4}{c}{15--40 parameters} \\
\bottomrule
\end{tabular}%
}
\end{table}

\Cref{tab:dataset_stats} summarizes the key statistics.
On average, each image undergoes $3.2 \pm 1.4$ distinct degradation operations, with 15--40 parameters recorded per image.
The stochastic nature of \pipeline ensures broad coverage of the degradation parameter space: JPEG quality factors span 25--85, sensor noise levels range over an order of magnitude, and all five motion blur types are represented.

\begin{figure}[t]
    \centering
    \includegraphics[width=\columnwidth]{figures/fig_dataset_samples.png}
    \caption{\textbf{Example LR--HR pairs from \dataset.}
    Left: HR ground truth ($1024 \times 1024$). Right: four LR variants ($256 \times 256$) produced by \pipeline at different difficulty levels, exhibiting realistic combinations of blur, noise, color shifts, and compression artifacts.}
    \label{fig:dataset_samples}
\end{figure}

% =================================================================================================
% 5. EXPERIMENTS
% =================================================================================================
\section{Experiments}
\label{sec:experiments}

We validate \dataset through a comprehensive evaluation of state-of-the-art SISR methods, demonstrating the degradation gap between standard benchmarks and our realistic degradation protocol.

\subsection{Setup}
\label{sec:exp_setup}

\paragraph{Evaluated methods.}
We select four representative models:
\textbf{DSCF-SR}~\cite{xiao2024dscf} (Transformer with dynamic sparse cross-attention),
\textbf{SeemoRe-B} and \textbf{SeemoRe-T}~\cite{zamfir2024seemore} (mixture-of-experts, Base and Tiny variants), and
\textbf{Real-ESRGAN}~\cite{wang2021real} (GAN-based blind SR with second-order degradation training).
The first three are trained under bicubic degradation on DIV2K/Flickr2K; Real-ESRGAN uses its own stochastic degradation pipeline.
All models use officially released $\times$4 weights; inference is on an NVIDIA Tesla P100 GPU.

\paragraph{Metrics.}
We employ PSNR and SSIM~\cite{wang2004image} (computed on the Y channel) as distortion metrics, and LPIPS~\cite{zhang2018unreasonable} as a learned perceptual quality metric.

\subsection{Overall Results}
\label{sec:exp_results}

\begin{table}[t]
\centering
\caption{\textbf{Overall performance on \dataset ($\times$4).}
Non-blind methods cluster in distortion metrics; the blind method (Real-ESRGAN) achieves lower PSNR but substantially better perceptual quality.}
\label{tab:overall_results}
\resizebox{\columnwidth}{!}{%
\begin{tabular}{@{}lcccc@{}}
\toprule
\textbf{Method} & \textbf{Training} & \textbf{PSNR$\uparrow$} & \textbf{SSIM$\uparrow$} & \textbf{LPIPS$\downarrow$} \\
\midrule
DSCF-SR~\cite{xiao2024dscf} & Bicubic & 23.82 & 0.674 & 0.556 \\
SeemoRe-B~\cite{zamfir2024seemore} & Bicubic & \textbf{23.86} & 0.673 & 0.561 \\
SeemoRe-T~\cite{zamfir2024seemore} & Bicubic & 23.85 & 0.673 & 0.560 \\
\midrule
Real-ESRGAN~\cite{wang2021real} & Blind deg. & 22.65 & \textbf{0.686} & \textbf{0.444} \\
\bottomrule
\end{tabular}%
}
\end{table}

\Cref{tab:overall_results} presents aggregate results across all 1{,}140 test pairs.
The best non-blind methods achieve only ${\sim}$23.9~dB PSNR on \dataset, compared to ${>}$30~dB on Set5 under bicubic degradation---a \textbf{6--8~dB gap} that directly quantifies the degradation gap motivating our work.
DSCF-SR, SeemoRe-B, and SeemoRe-T converge to nearly identical performance (within 0.04~dB), suggesting that under complex degradations the performance ceiling is determined by the degradation mismatch rather than architectural capacity.

Real-ESRGAN presents a striking counterpoint: while achieving the lowest PSNR ($-$1.21~dB below non-blind methods), it obtains the best SSIM (0.686) and substantially better LPIPS (0.444 vs.\ ${>}$0.555, a 20\% relative improvement).
This perception--distortion divergence~\cite{blau2018perception} underscores the need for multi-metric evaluation under complex degradations.

\subsection{Category-Wise Analysis}
\label{sec:exp_category}

\begin{table*}[t]
\centering
\caption{\textbf{Category-wise performance on \dataset ($\times$4).}
Best per metric and category in \textbf{bold}.
Urban architecture is most challenging; Real-ESRGAN leads LPIPS across all categories.}
\label{tab:category_results}
\resizebox{\textwidth}{!}{%
\begin{tabular}{@{}l|ccc|ccc|ccc|ccc@{}}
\toprule
& \multicolumn{3}{c|}{\textbf{Urban Architecture}} & \multicolumn{3}{c|}{\textbf{Natural Landscapes}} & \multicolumn{3}{c|}{\textbf{Portraits \& People}} & \multicolumn{3}{c}{\textbf{Macro Objects}} \\
\cmidrule(l){2-4} \cmidrule(l){5-7} \cmidrule(l){8-10} \cmidrule(l){11-13}
\textbf{Method} & PSNR$\uparrow$ & SSIM$\uparrow$ & LPIPS$\downarrow$ & PSNR$\uparrow$ & SSIM$\uparrow$ & LPIPS$\downarrow$ & PSNR$\uparrow$ & SSIM$\uparrow$ & LPIPS$\downarrow$ & PSNR$\uparrow$ & SSIM$\uparrow$ & LPIPS$\downarrow$ \\
\midrule
DSCF-SR & 23.26 & 0.692 & 0.505 & 23.12 & 0.676 & 0.539 & \textbf{24.55} & 0.658 & 0.602 & 24.44 & 0.669 & 0.581 \\
SeemoRe-B & \textbf{23.35} & 0.693 & 0.511 & \textbf{23.14} & 0.674 & 0.545 & \textbf{24.58} & 0.657 & 0.607 & \textbf{24.48} & 0.667 & 0.584 \\
SeemoRe-T & 23.34 & 0.692 & 0.510 & 23.13 & 0.674 & 0.544 & 24.57 & 0.657 & 0.606 & 24.48 & 0.667 & 0.584 \\
\midrule
Real-ESRGAN & 21.89 & \textbf{0.702} & \textbf{0.383} & 22.17 & \textbf{0.701} & \textbf{0.408} & 23.16 & \textbf{0.643} & \textbf{0.532} & 23.43 & \textbf{0.694} & \textbf{0.463} \\
\bottomrule
\end{tabular}%
}
\end{table*}

\Cref{tab:category_results} reveals content-dependent vulnerability patterns.
\textbf{Urban architecture} is consistently the most challenging category (${\sim}$23.3~dB for non-blind methods vs.\ ${\sim}$24.5~dB for portraits/macro), as geometric structures interact adversely with blur and compression artifacts.
\textbf{Macro objects} are comparatively resilient due to textural regularity and self-similar structures.
Real-ESRGAN's LPIPS advantage is universal (0.119--0.131 improvement), largest for natural landscapes where adversarially generated organic textures are perceptually more convincing.
In portraits, the PSNR gap narrows ($-$1.39~dB) while the LPIPS gap widens, highlighting the acute perception--distortion tradeoff for facial content.

\subsection{Standard Benchmarks vs.\ \dataset}
\label{sec:exp_gap_analysis}

\begin{table}[t]
\centering
\caption{\textbf{Performance gap: bicubic benchmarks vs.\ \dataset.}
The 6--8~dB drop reveals a severe generalization gap.}
\label{tab:gap_analysis}
\resizebox{\columnwidth}{!}{%
\begin{tabular}{@{}lccccc@{}}
\toprule
\textbf{Method} & \textbf{Set5} & \textbf{Set14} & \textbf{BSD100} & \textbf{Urban100} & \textbf{\dataset} \\
\midrule
Non-blind SOTA & $\sim$33 & $\sim$29 & $\sim$28 & $\sim$27 & $\sim$23.9 \\
Real-ESRGAN & $\sim$28 & $\sim$25 & $\sim$25 & $\sim$24 & 22.7 \\
\midrule
$\Delta$ (drop) & -- & -- & -- & -- & 6--8 dB \\
\bottomrule
\end{tabular}%
}
\end{table}

\Cref{tab:gap_analysis} contextualizes the degradation gap.
Even Real-ESRGAN, explicitly trained for blind SR, drops ${\sim}$5~dB relative to Urban100, indicating that degradation-aware training has substantial room for improvement against the diversity of \pipeline.
The high per-image variance (5.7--7.1~dB PSNR standard deviation) further highlights that model performance on realistic data is highly image-dependent, motivating per-image metadata for fine-grained analysis.

\subsection{Qualitative Results}
\label{sec:exp_qualitative}

\begin{figure*}[t]
    \centering
    \includegraphics[width=0.92\textwidth]{figures/fig_qualitative_comparison.png}
    \caption{\textbf{Qualitative comparison on worst-case samples from \dataset.}
    \textbf{Top (a--b):} Non-blind methods fail catastrophically on realistic degradations---severe illumination causes darkened outputs (PSNR ${<}$9~dB), and texture ringing emerges on fine structures.
    \textbf{Bottom (c--d):} Real-ESRGAN struggles with extreme color/exposure shifts.
    These failures are invisible in standard bicubic benchmarks.}
    \label{fig:qualitative_comparison}
\end{figure*}

\Cref{fig:qualitative_comparison} shows representative failure cases.
Non-blind methods exhibit characteristic artifacts: noise amplification, color fringing from chromatic aberration, and ringing where the assumed bicubic kernel differs from the actual degradation.
Real-ESRGAN produces perceptually more coherent outputs but occasionally hallucinates textures not present in the ground truth.

\subsection{Toward Degradation-Aware SR}
\label{sec:exp_dast}

The per-image degradation metadata in \dataset enables degradation-conditioned SR.
As a proof of concept, we describe a baseline Degradation-Aware Swin Transformer (DAST) that predicts a degradation vector from the LR input and conditions the backbone via SFT layers~\cite{wang2018recovering}: $\mathbf{F}' = (\bm{\gamma}(\mathbf{d}) + \mathbf{1}) \odot \mathbf{F} + \bm{\beta}(\mathbf{d})$.
The estimator is supervised using ground-truth metadata, enabling precise characterization without iterative refinement---a capability impossible to develop with existing benchmarks.

% =================================================================================================
% 6. CONCLUSION
% =================================================================================================
\section{Conclusion}
\label{sec:conclusion}

We have presented \dataset, a benchmark for evaluating SISR under realistic, physically-grounded degradations, accompanied by \pipeline---a multi-stage degradation pipeline modeling the image formation process from optical aberrations to lossy compression.
Our benchmarking of four SOTA methods reveals a 6--8~dB PSNR gap between standard bicubic benchmarks and \dataset, with non-blind methods converging to similar performance regardless of architecture, and a striking perception--distortion divergence for the blind method Real-ESRGAN.

The per-image degradation metadata---unique to \dataset---opens several research directions: degradation-conditioned models that adapt to input-specific degradation profiles (\cf \cref{sec:exp_dast}), stratified robustness evaluation by degradation type and severity, and curriculum learning with calibrated difficulty levels.
The modular design of \pipeline allows extension with new degradation types or adaptation to other restoration tasks.
We will publicly release the complete dataset, pipeline code, and evaluation scripts to facilitate reproducible research.


\FloatBarrier
{
    \small
    \bibliographystyle{ieeenat_fullname}
    \bibliography{main}
}

\end{document}
